\documentclass[10pt,a4paper]{amsart}
\usepackage[margin=19mm]{geometry}
\usepackage{amsmath,amssymb,mathtools}
\usepackage[scr=rsfs,cal=euler]{mathalfa}
\usepackage{xfrac}
\usepackage[unicode,hidelinks,bookmarks=false]{hyperref}
\newcommand{\ie}{i.e.\ }
\newcommand{\cf}{cf.\ }
\newcommand{\resp}{respectively\ }
\newcommand{\Subsection}{\subsection}
\providecommand{\nobreakdash}{\nobreak\hbox{-}}
\setlength{\emergencystretch}{2em}
\multlinegap0pt
\usepackage{mathrsfs}
\usepackage{mathtools}
\usepackage{enumitem}
\usepackage{tikz}
\usepackage{tikz-cd}
\usepackage{xy}
\newtheorem{thm}{Theorem}[section]
\newtheorem{conjecture}[thm]{Conjecture}
\newtheorem{cor}[thm]{Corollary}
\newtheorem{conj}[thm]{Conjecture}
\newtheorem{lemma}[thm]{Lemma}
\newtheorem{lem}[thm]{Lemma}
\newtheorem{prop}[thm]{Proposition}
\newtheorem{problem}[thm]{Problem}
\newtheorem{example}[thm]{Example}
\newtheorem{ex}[thm]{Example}
\newtheorem{defn}[thm]{Definition}
\newtheorem{nota}[thm]{Notation}
\newtheorem{rmk}[thm]{Remark}
\newtheorem{ques}[thm]{Question}
\setlength{\marginparwidth}{2cm}
\numberwithin{equation}{section}
\usepackage{todonotes}
\newcommand\du[1]{\todo[inline,color=green!40]{#1}} %Du Jiabin inline
\newcommand\hu[1]{\todo[inline,color=yellow!40]{#1}} %Hu Yong inline
\newcommand{\alb}{\mathrm{alb}}
\newcommand{\Alb}{\mathrm{Alb}}
\newcommand{\alg}{\mathrm{alg}}
\newcommand{\CH}{\mathrm{CH}}
\newcommand{\ev}{\mathrm{ev}}
\newcommand{\hhom}{\mathrm{hom}}
\newcommand{\id}{\mathrm{id}}
\newcommand{\im}{\mathrm{Im}}
\newcommand{\Jac}{\mathrm{Jac}}
\newcommand{\Pic}{\mathrm{Pic}}
\newcommand{\rk}{\mathrm{rank}}
\newcommand{\supp}{\mathrm{supp}}
\newcommand{\Sym}{\mathrm{Sym}}
\newcommand{\tr}{\mathrm{tr}}
\newcommand{\Aut}{\mathrm{Aut}}
\newcommand{\Albanese}{\mathrm{Albanese}}
\newcommand{\Kimura}{\mathrm{Kimura}}
\newcommand{\MW}{\mathrm{MW}}
\newcommand{\End}{\mathrm{End}}
\newcommand{\exc}{\mathrm{Ex}}
\newcommand{\cod}{\mathrm{codim}}
\newcommand{\ord}{\mathrm{ord}}
\newcommand{\Mot}{\mathrm{Mot}}
\newcommand{\CHM}{\mathrm{\CHM}}
\newcommand{\Corr}{\mathrm{Corr}}
\newcommand{\HHom}{\mathrm{Hom}}
\newcommand{\cl}{\mathrm{cl}}
\newcommand{\CC}{\mathbb{C}}
\newcommand{\QQ}{\mathbb{Q}}
\newcommand{\RR}{\mathbb{R}}
\newcommand{\ZZ}{\mathbb{Z}}
\newcommand{\PP}{\mathbb{P}}
\newcommand{\sB}{\mathcal{B}}
\newcommand{\sE}{\mathcal{E}}
\newcommand{\sF}{\mathcal{F}}
\newcommand{\sH}{\mathcal{H}}
\newcommand{\sM}{\mathcal{M}}
\newcommand{\sO}{\mathcal{O}}
\newcommand{\gh}{\mathfrak{h}}
\newcommand{\tA}{\widetilde{A}}
\newcommand{\tB}{\widetilde{B}}
\newcommand{\tC}{\widetilde{C}}
\newcommand{\tD}{\widetilde{D}}
\newcommand{\tE}{\widetilde{E}}
\newcommand{\tF}{\widetilde{F}}
\newcommand{\tG}{\widetilde{G}}
\newcommand{\tH}{\widetilde{H}}
\newcommand{\tI}{\widetilde{I}}
\newcommand{\tJ}{\widetilde{J}}
\newcommand{\tK}{\widetilde{K}}
\newcommand{\tL}{\widetilde{L}}
\newcommand{\tM}{\widetilde{M}}
\newcommand{\tN}{\widetilde{N}}
\newcommand{\tO}{\widetilde{O}}
\newcommand{\tP}{\widetilde{P}}
\newcommand{\tQ}{\widetilde{Q}}
\newcommand{\tR}{\widetilde{R}}
\newcommand{\tS}{\widetilde{S}}
\newcommand{\tT}{\widetilde{T}}
\newcommand{\tU}{\widetilde{U}}
\newcommand{\tV}{\widetilde{V}}
\newcommand{\tW}{\widetilde{W}}
\newcommand{\tX}{\widetilde{X}}
\newcommand{\tY}{\widetilde{Y}}
\newcommand{\tZ}{\widetilde{Z}}
\newcommand{\OO}{\mathcal{O}}
\newcommand{\bo}{\omega}
\def\red{\color[rgb]{1,0,0}}
\def\blu{\color[rgb]{0,0,1}}
\begin{document}
\title[On the canonical degree of a Gorenstein minimal threefold of]{On the canonical degree of a Gorenstein minimal threefold of general type}
\author{Dongchuan Road 800; Jiabin Du; Yong Hu}
\date{}
\maketitle
\noindent\emph{Source and licence.} On the canonical degree of a Gorenstein minimal threefold of general type. Version arXiv:2606.31170v1 (CC BY 4.0). This material is distributed under the Creative Commons Attribution 4.0 International licence, http://creativecommons.org/licenses/by/4.0/, as recorded in the arXiv OAI licence metadata for this exact version and shown on the abstract page https://arxiv.org/abs/2606.31170v1. Full rights notice: licenses/arxiv-2606.31170-CC-BY-4.0.txt.\par\smallskip
\noindent\textit{Editorial scope.} Counted unit: the original \texttt{cor} environment \texttt{MY} at source lines 257--263 together with its complete proof at lines 264--274; the delivered document keeps source lines 228--275 so that every referenced label and every citation resolves inside it. Native editable source is the paper's own concatenated TeX; the AMS wrapper is editorial formatting only. No publisher class, upstream script or unrelated artwork is redistributed.
\par\smallskip\noindent\textit{Reference note.} This article ships no compiled bibliography (biblatex); the entries delivered here are composed from the author's own BibTeX (.bib) fields, with no field invented and no entry dropped.
\section{Preliminaries}



\subsection{Image sheaf of the evaluation map}
Let $C$ be a smooth  curve of genus $g$, and let $E$ be a vector bundle of rank $r$ on $C$. Consider the evaluation map
\[
\ev\colon H^0(E)\otimes\sO_C\to E
\]
Denote by $E_{\ev}$ the image sheaf of $\ev$, which is a vector bundle, say of rank $l$.
\begin{thm}\cite[Proposition~page 476]{Xia86}\label{Xiao_Eev}
    Suppose that $E\otimes\omega_C$ is nef, then
    \[
    (g-1)(r-l)\leq l.
    \]
    Moreover, when the equality holds, $E$ is a semi-stable bundle of degree $(2g-2)r$.
\end{thm}

\subsection{Miyaoka-Yau inequality}

For minimal varieties of general type, Greb, Kebekus, Peternell and Taji (see \cite[Theorem 1.1]{GKPT19} and \cite[line 12 in \S 1.1, page 1488]{GKPT19}) proved the following $\QQ$-Miyaoka-Yau inequality:

\begin{thm}[$\mathbb{Q}$-Miyaoka-Yau inequality]\label{thm: MY1}
Let $X$ be an $n$-dimensional, minimal variety of general type. Then,
\begin{equation*}
\bigl(2(n+1)\cdot c_2(X) - n\cdot c_1(X)^2\bigr)\cdot [K_X]^{n-2} \geq 0.
\end{equation*}
\end{thm}
In particular, we have

\begin{cor}
    Let $X$ be a Gorenstein minimal $3$-fold of general type, then 
    \begin{equation}\label{MY}
        K_X^3\leq 64\chi(\omega_X).
    \end{equation}
   
\end{cor}

\begin{proof}
    By Theorem \ref{thm: MY1}, we have 
    $$
    8(K_X\cdot c_2(X))\geq 3K_X^3.
    $$
    By \cite[Corollary~10.3]{Rei87}, we have $(K_X\cdot c_2(X))=24\chi(\omega_X)$. It follows easily that 
    $$
    K_X^3\leq 64\chi(\omega_X).
    $$
    The proof is completed.
\end{proof}


\begin{thebibliography}{99}
\bibitem[]{GKPT19} Greb, Daniel; Kebekus, Stefan; Peternell, Thomas; Taji, Behrouz. The {M}iyaoka-{Y}au inequality and uniformisation of canonical models. Ann. Sci. \'Ec. Norm. Sup\'er. (4) 52 (2019) 1487--1535
\bibitem[]{Rei87} Reid, Miles. Young person's guide to canonical singularities. Algebraic geometry, {B}owdoin, 1985 ({B}runswick, {M}aine, 1985) 46, Part 1 (1987) 345--414
\bibitem[]{Xia86} Xiao, Gang. Algebraic surfaces with high canonical degree. Math. Ann. 274 (1986) 473--483
\end{thebibliography}
\end{document}
